\chapter{Shigefumi Mori}\label{ch1990shigefumimori:shigefumi-mori-fields-medal-1990-a-comprehensive-report}

\begin{center}\rule{0.5\linewidth}{0.5pt}\end{center}

\textbf{Prepared:} July 2026 \textbf{Subject:} Shigefumi Mori (b.~February 17, 1951, Nagoya, Japan)
\textbf{Position:} Professor, Kyoto University (since 1990)
\textbf{Doctoral Advisor:} Masayoshi Nagata (Kyoto University)
\textbf{Fields Medal Citation:} \emph{``For his work on the classification of algebraic three-folds and the construction of new algebraic three-folds.''}

\begin{center}\rule{0.5\linewidth}{0.5pt}\end{center}

\begin{prereqs}
Algebraic geometry (divisor theory, Mori theory, birational geometry),
minimal model program, cone theorem, deformation theory of curves.
\end{prereqs}

\begin{objectives}
State Mori's cone theorem and the bend-and-break lemma;
explain the minimal model program for three-folds;
describe the classification of Fano three-folds;
understand the existence of flips and the log minimal model program.
\end{objectives}

\section{Problem History}

The classification of algebraic varieties---specifically, finding a canonical
representative in each birational equivalence class---is one of the central
problems in algebraic geometry. For surfaces (dimension 2), the classification
was completed by Enriques, Kodaira, and Zariski in the early 20th century.
Every surface is birationally equivalent to either a minimal surface (with no
$(-1)$-curves) or a ruled surface. The minimal model is unique up to
isomorphism.

For varieties of dimension $\geq 3$, the minimal model program was expected to
generalize the surface case, but the obstacles were much more severe. In
dimension 3, Mori \cite{mori1982} introduced new techniques to construct minimal models,
fundamentally changing the understanding of birational geometry.

The central problem is to find, for every smooth projective variety $X$, either
a minimal model (a smooth variety with nef canonical divisor $K_X$) or a Mori
fiber space (a variety fibred by Fano varieties). The existence of such a
canonical representative would provide a complete classification of algebraic
varieties up to birational equivalence.

The minimal model program proceeds through a sequence of birational transformations.
Starting from a smooth projective variety $X$, one seeks to contract curves $C$
with negative intersection $K_X \cdot C < 0$ while preserving the "goodness" of
the variety. The process should terminate in either:
\begin{itemize}
    \item A \textbf{minimal model}: a variety $X_{\min}$ with nef canonical divisor
    $K_{X_{\min}}$, where no further contractions of curves are possible; or
    \item A \textbf{Mori fiber space}: a fibration $f: X \to Y$ with $\dim Y < \dim X$
    whose general fiber is a Fano variety (i.e., has ample anticanonical divisor).
\end{itemize}

Mori's work was motivated by the Hartshorne conjecture, which predicted that
every smooth subvariety of $\mathbb{P}^n$ of codimension $\leq n/2$ is a complete
intersection. Mori proved the conjecture for codimension $\leq n/2$ in 1975,
using the theory of varieties of small degree. This result showed that the theory
of extremal rays---the extreme rays of the cone of curves---could be used to
study the structure of algebraic varieties.

The theory of varieties of small degree begins with a fundamental question:
given a smooth projective variety $X$ of dimension $n$ embedded in $\mathbb{P}^N$
by a very ample line bundle $L$, when does the degree $d = L^n$ impose strong
restrictions on the structure of $X$? Mori showed that small degree forces $X$
to be highly special.

For a smooth projective variety $X$ of dimension $n$, the \textbf{degree} of $X$
with respect to an embedding by a very ample line bundle $L$ is defined as
$d = L^n = \deg(X)$. The varieties of small degree are those satisfying
$d \leq 2n - 1$. Mori's classification theorem states:

\begin{theorem}[Mori's Classification of Varieties of Small Degree]
Let $X$ be a smooth projective variety of dimension $n$ with a very ample
line bundle $L$ such that $L^n \leq 2n - 1$. Then $(X, L)$ is one of the
following:
\begin{enumerate}
    \item $(\mathbb{P}^n, \mathcal{O}(1))$: projective space with hyperplane bundle;
    \item $(Q^n, \mathcal{O}(1))$: a quadric hypersurface in $\mathbb{P}^{n+1}$;
    \item A Del Pezzo surface: a blow-up of $\mathbb{P}^2$ at $k \leq 8$ points;
    \item A scroll over a curve: a projective bundle $\mathbb{P}(E)$ over a curve $C$
    of genus $g$, with $L$ a relatively ample generator.
\end{enumerate}
\end{theorem}

This classification provided the first major application of the theory of
extremal rays to the structure of projective varieties. The key insight is that
small degree implies the existence of a covering family of lines (rational curves
of minimal degree), which can be used to deform and construct new varieties.

The bend-and-break lemma, introduced by Mori, is a deformation-theoretic tool
that shows: if $X$ is a smooth projective variety and $C$ is a rational curve
of small degree, then $C$ can be deformed to break into a sum of rational curves.
This lemma is the key technical tool in Mori's theory and has many applications
beyond the minimal model program.

The \textbf{bend-and-break lemma} can be stated as follows. Let $X$ be a smooth
projective variety of dimension $n$, and let $l$ be a line in $\mathbb{P}^N$
interacting with $X$ in a suitable way. Then there exists a family of rational
curves in $X$ covering $X$, and as the parameter varies, the curve can bend
until it breaks into a union of rational curves. This breaking phenomenon is
crucial for constructing the extremal rays of the cone of curves.

The cone of curves $\overline{NE}(X)$, the closure of the cone generated by
classes of effective curves in $X$, is a fundamental object in Mori's theory.
The cone theorem describes the structure of $\overline{NE}(X)$ in terms of the
canonical divisor $K_X$. The theorem shows that $\overline{NE}(X)$ is generated
by a finite number of extremal rays, each of which corresponds to a birational
contraction or fiber space structure.

Mori's work on three-folds in the 1980s led to the classification of smooth
Fano three-folds---three-folds with ample anticanonical divisor. Fano three-folds
are the building blocks of the minimal model program: every algebraic variety is
birationally equivalent to a Mori fiber space, which is a Fano variety fibred
over a lower-dimensional base. The classification of Fano three-folds proceeds
by analyzing the index $r$, defined as the largest integer dividing $-K_X$ in
the Picard group. For index $r \geq 2$, the Fano three-folds are
well-understood:
\begin{itemize}
    \item $r = 4$: $\mathbb{P}^3$;
    \item $r = 3$: quadric hypersurface $Q \subset \mathbb{P}^4$;
    \item $r = 2$: Del Pezzo three-fold (complete intersection of two quadrics
    or a cubic in weighted projective space);
    \item $r = 1$: 17 deformation families with various properties.
\end{itemize}

The \textbf{index} of a Fano manifold $X$ is defined as the largest positive
integer $r$ such that $-K_X = rH$ for some ample divisor $H$ on $X$. The
integer $r$ is called the \textbf{Fano index} and satisfies $1 \leq r \leq
n+1$, where $n = \dim X$. For $r = n+1$, $X \cong \mathbb{P}^n$. For $r = n$,
$X$ is a quadric hypersurface. These results are the starting point for the
classification of Fano varieties.

The difficulty in extending the MMP to higher dimensions lies in the phenomenon
of flips. A flip is a birational transformation that replaces one extremal ray
with another, producing a new variety with the same Picard number but different
geometry. In dimension 3, Mori proved that flips exist and terminate, but in
higher dimensions, the existence and termination of flips are major open
problems.

The theory of \textbf{log terminal singularities} plays a crucial role in the MMP.
A variety has log terminal singularities if the discrepancies of all exceptional
divisors are strictly greater than $-1$. The existence of minimal models for
varieties with log terminal singularities is a key result of the MMP, and the
classification of log terminal singularities in dimension 3 has been completed
by several researchers.

The bend-and-break technique has applications beyond the minimal model program.
It has been used to study the existence of rational curves on Fano manifolds,
the connectedness of the space of morphisms between algebraic varieties, and the
structure of the Chow ring of rationally connected varieties. The technique has
also been generalized to the setting of K\"ahler manifolds, where it provides
tools for studying the geometry of holomorphic curves.

The MMP for three-folds has been refined to account for singularities. The log
MMP, developed by Kawamata and others, deals with pairs $(X, \Delta)$ where
$\Delta$ is a $\mathbb{Q}$-divisor with coefficients in $[0,1]$. The log
canonical ring of a pair $(X, \Delta)$ is finitely generated when $K_X + \Delta$
is big, and the log minimal model is constructed using the same techniques as
in the smooth case. The classification of log canonical singularities in
dimension 3, completed by Kawamata and others, provides the foundation for the
log MMP.

\section{How Was It Solved}

The central result of Mori's theory is the \textbf{cone theorem}, which describes
the structure of the cone of curves $\overline{NE}(X)$ for a smooth projective
variety $X$ of dimension $n$. The theorem provides a decomposition of
$\overline{NE}(X)$ into a sum of extremal rays and the nef cone.

\begin{theorem}[Mori's Cone Theorem]
Let $X$ be a smooth projective variety of dimension $n$ over an algebraically
closed field of characteristic zero. Then:
\[
\overline{NE}(X) = \overline{NE}(X)_{K_X \geq 0} + \sum_{i} R_i
\]
where $R_i$ are the extremal rays with $K_X \cdot R_i < 0$, and
$\overline{NE}(X)_{K_X \geq 0}$ is the cone of curves $C$ satisfying
$K_X \cdot C \geq 0$. The sum is countable, and each $R_i$ is generated by a
rational curve $C_i$ with $-n \leq K_X \cdot C_i < 0$.
\end{theorem}

The cone theorem has several important consequences. First, it implies that the
nef cone $\text{Nef}(X)$, which is the dual cone to $\overline{NE}(X)$, has a
describable structure: its boundary is composed of hyperplanes orthogonal to
the extremal rays. Second, it provides the foundation for the minimal model
program by identifying the curves that should be contracted.

The proof of the cone theorem uses the bend-and-break lemma. Given a curve
$C$ with small degree, one can deform $C$ by bending it until it breaks into a
sum of rational curves. The bend-and-break deformation is achieved by finding a
family of curves parametrized by a quasi-projective variety, and using the fact
that rational curves have non-trivial automorphisms to break the curve.

The key technical tool is the \textbf{moving lemma}: one can find a moving
rational curve---a rational curve in a continuous family---to bend and break.
The moving lemma is proved using the theory of low-degree curves on varieties,
based on Mori's classification of varieties of small degree.

The existence of rational curves of low degree follows from Mori's theorem on
varieties of small degree: if $L$ is an ample line bundle on $X$ with
$L^n \leq 2\dim X + 1$ (where $n = \dim X$), then $X$ is a linear space, a
quadric, or a projective bundle. This classification provides the rational
curves needed for bend-and-break.

The \textbf{contraction theorem} states: every extremal ray $R$ of
$\overline{NE}(X)$ is contractible. That is, there exists a morphism
$\text{cont}_R: X \to Y$ to a normal projective variety $Y$ such that
$\text{cont}_R$ contracts exactly those curves whose classes lie in $R$. The
proof uses the theory of polarized pairs and the classification of Fano
varieties.

More precisely, given an extremal ray $R$, the contraction theorem guarantees
the existence of a projective morphism $\phi: X \to Y$ with connected fibers
such that a curve $C$ is contracted by $\phi$ if and only if $[C] \in R$. The
target $Y$ is a normal projective variety, and the relative Picard number
satisfies $\rho(X/Y) = 1$.

For three-folds, the MMP terminates because the dimension decreases at each
birational step: a birational contraction of a three-fold produces a two-fold,
which is already classified. The general MMP for higher-dimensional varieties
is still open, but the theory of weak log canonical models and K\"ahler--Einstein
metrics has led to significant progress.

The bend-and-break lemma also implies the existence of rational curves on
rationally connected varieties: if $X$ is rationally connected, then through
any two points there exists a rational curve. This provides a geometric
characterization of varieties that are birationally equivalent to Fano varieties.

The proof of the cone theorem for three-folds uses the following key inequality.
For an ample divisor $H$ and an extremal ray $R$ with $K_X \cdot R < 0$,
there exists a positive integer $a$ and a rational number $b$ such that for
every curve $C \in R$:
\[
H \cdot C \leq a(-K_X \cdot C) + b
\]
This bound controls the geometry of curves in the extremal ray and is essential
for the construction of the contraction morphism.

\subsection*{Classification of Fano Three-Folds}

The classification of smooth Fano three-folds is one of Mori's greatest
achievements. The classification proceeds by analyzing the index $r$ and the
Picard number $\rho$. The main results are:

\begin{theorem}[Classification of Fano Three-Folds]
Let $X$ be a smooth Fano three-fold. Then:
\begin{enumerate}
    \item If $\rho(X) = 1$, then $X$ belongs to one of 17 deformation families,
    distinguished by the degree $(-K_X)^3$ and the index $r$.
    \item If $\rho(X) > 1$, then $X$ can be obtained from a Fano three-fold
    of Picard number 1 by a sequence of blow-ups of smooth curves.
\end{enumerate}
\end{theorem}

The 17 families of Fano three-folds with Picard number 1 are distinguished by
their index $r$ and the degree $d = (-K_X)^3$. For index $r \geq 2$, the
classification is complete and explicit. The families with $r = 1$ are more
complicated and were classified by Mori and Mukai.

The degree $d = (-K_X)^3$ satisfies $1 \leq d \leq 64$, with $d = 64$ giving
$\mathbb{P}^3$ and $d = 54$ giving the quadric $Q \subset \mathbb{P}^4$. The
possible values of $d$ depend on the index $r$, and the classification involves
a careful analysis of the anticanonical linear system $|-K_X|$.

For $r = 1$, the degree $d$ can take values $1 \leq d \leq 64$ with various
constraints. The case $d = 1$ is particularly interesting: a Fano three-fold
of index 1 and degree 1 is a double cover of a weighted projective space
$\mathbb{P}(1,1,1,2,3)$ branched along a surface of degree 6.

\subsection*{Existence of Flips}

A major breakthrough in Mori's work was the proof of the existence of flips.
A flip is a birational transformation that replaces a curve with negative
intersection with the canonical divisor by another curve with positive
intersection. Flips are necessary for the MMP to proceed beyond the first step.

\begin{theorem}[Mori's Flip Theorem]
Let $X$ be a smooth three-fold and let $R$ be an extremal ray with $K_X \cdot
R < 0$. Then the flip of $R$ exists: there exists a birational morphism
$\phi^+: X^+ \to Y$ such that $K_{X^+}$ is ample over $Y$ and $\phi^+$ is an
isomorphism in codimension 1.
\end{theorem}

The proof of the flip theorem is highly technical and involves several steps:
\begin{enumerate}
    \item Constructing the flip as a limit of a sequence of blow-ups;
    \item Using the bend-and-break lemma to produce rational curves on the flip;
    \item Applying the classification of varieties of small degree to control
    the geometry;
    \item Proving that the flip is projective and has the required properties.
\end{enumerate}

The existence of flips in dimension 3 opened the door to the full minimal model
program. However, the existence of flips in higher dimensions remained open for
many years. The general case was finally resolved by Birkar, Cascini, Hacon,
and McKernan in 2010.

\subsection*{Termination of Flips}

In dimension 3, Mori also proved that the MMP terminates: a sequence of flips
must eventually stop. This is a consequence of the fact that the Picard number
decreases at each step, and there are only finitely many possible Picard numbers.

The termination of flips in higher dimensions is still open in general, but
significant progress has been made. For varieties of dimension 4, the termination
of flips has been proved by Birkar and others. The general case remains one of
the central open problems in birational geometry.

\section{Impact of This Discovery}

Mori's theory provides the foundation for modern birational geometry. The cone
theorem and bend-and-break technique have been generalized to many settings:
singular varieties, K\"ahler manifolds, and non-archimedean geometry. The
minimal model program is the central organizing principle for the classification
of algebraic varieties.

The classification of Fano three-folds, completed by Iskovskikh, Mori, and
Mukai, shows that there are exactly 105 deformation families of smooth Fano
three-folds. This classification is one of the great achievements of 20th
century algebraic geometry and provides a complete picture of the building
blocks of three-fold geometry.

The MMP has been extended to the setting of log geometry. The log MMP deals
with pairs $(X, \Delta)$, where $\Delta$ is a divisor with coefficients in $[0,1]$.
The log cone theorem, proved by Kawamata and others, extends Mori's cone theorem
to the log setting, providing a more flexible framework for birational geometry.

The relationship between the MMP and K\"ahler--Einstein metrics, discovered by
Tian, Donaldson, and others, connects birational geometry to differential
geometry. The existence of a K\"ahler--Einstein metric on a Fano manifold is
equivalent to the existence of a ``K\"ahler--Ricci soliton,'' and the stability
of the manifold in the sense of GIT is necessary and sufficient for the
existence of such a metric. This is the Yau--Tian--Donaldson conjecture, proved
by Chen--Donaldson--Sun and Tian.

Mori's work also influenced the study of rational curves on algebraic varieties.
The enumerative geometry of rational curves---counting rational curves of given
degree on a variety---is now a major area of research, connected to
Gromov--Witten invariants and mirror symmetry. The bend-and-break lemma
provides the foundation for understanding how rational curves can be deformed
and degenerated.

The MMP for three-folds has been implemented in computer algebra systems by
the Birational Geometry Project. The explicit classification of all birational
types of three-folds, and the construction of explicit minimal models, is now
feasible for many specific three-folds. This computational aspect has made the
MMP a practical tool for algebraic geometers.

The notion of a ``minimal model'' has been refined to account for the presence
of singularities. The canonical model of a variety of general type, which has
mild singularities (log terminal singularities), provides a unique birational
representative. The existence of canonical models in all dimensions, proved by
Birkar, Cascini, Hacon, and McKernan, is a major achievement of the MMP.

\subsection*{Rational Curves and Mori Space}

Mori's bend-and-break lemma led to the introduction of the \textbf{Mori space}
(or Kontsevich moduli space) $\overline{M}_{0,n}(X, \beta)$, which parametrizes
stable maps from genus 0 curves to $X$ with degree $\beta$. This space is a
fundamental tool in enumerative geometry and Gromov--Witten theory.

The Mori space has the following properties:
\begin{itemize}
    \item It is a proper algebraic space (and a scheme in many cases);
    \item It carries evaluation maps $\text{ev}_i: \overline{M}_{0,n}(X, \beta) \to X$
    sending a stable map to the image of the $i$-th marked point;
    \item The Gromov--Witten invariants are defined by integrating cohomology
    classes over the virtual fundamental class of $\overline{M}_{0,n}(X, \beta)$.
\end{itemize}

These invariants count the number of rational curves of degree $\beta$ passing
through specified cycles in $X$, and they play a central role in mirror symmetry
and string theory.

\subsection*{Rationally Connected Varieties}

Mori's work led to the theory of rationally connected varieties. A smooth
projective variety $X$ is \textbf{rationally connected} if any two general
points can be joined by a rational curve. This property is birational: if $X$
and $Y$ are birational, then $X$ is rationally connected if and only if $Y$
is.

Key results in the theory of rationally connected varieties include:
\begin{itemize}
    \item A Fano variety is rationally connected (proven using bend-and-break);
    \item A rationally connected variety has trivial rational cohomology in
    positive degrees;
    \item Rational connectedness is preserved under deformation and fibration.
\end{itemize}

The study of rationally connected varieties has applications to the
generalized Hodge conjecture, the Bloch--Kas--Lieberman conjecture, and the
study of algebraic cycles.

\subsection*{Applications to String Theory}

Mori's work has profound connections to string theory. The minimal model program
provides the geometric framework for understanding the compactification of
extra dimensions in string theory. Calabi--Yau three-folds, which are central
to string compactification, are closely related to Fano three-folds through
duality.

The mirror symmetry conjecture, which relates the enumerative geometry of
rational curves on a Calabi--Yau three-fold to the variation of Hodge structure
on its mirror, has been verified in many cases using techniques from the MMP.
The prediction of the number of rational curves of a given degree on a quintic
three-fold in $\mathbb{P}^4$ was one of the most famous successes of mirror
symmetry.

\section{Modern Views and Improvements}

The MMP for higher-dimensional varieties has been developed by many researchers.
Kawamata, Shokurov, Birkar, Cascini, Hacon, and McKernan have proved the
existence of minimal models for varieties of general type in all dimensions.
The weak log canonical approximation theorem and the termination of flips are
major open problems.

The connection to K\"ahler geometry has led to new approaches to the MMP. The
existence of K\"ahler--Einstein metrics on Fano manifolds, proved by
Chen--Donaldson--Sun and Tian, shows that the MMP is closely related to the
existence of canonical metrics. The Yau--Tian--Donaldson conjecture, now a
theorem, provides a differential-geometric characterization of the MMP.

The MMP has been extended to the non-archimedean setting. Berkovich spaces and
adic spaces provide a framework for studying birational geometry over
non-archimedean fields. The MMP for these spaces, developed by Musta\c{t}\v{a},
Nicaise, and others, has applications to arithmetic geometry and number theory.

The relationship between the MMP and tropical geometry is a recent development.
Tropical varieties, which are piecewise-linear objects associated to algebraic
varieties, encode the combinatorics of the MMP. The tropical MMP, studied by
Gross, Siebert, and others, provides a combinatorial framework for birational
geometry.

Mori himself has continued to make major contributions to algebraic geometry.
His work on the classification of algebraic varieties, the theory of Fano
manifolds, and the study of rational curves has influenced generations of
mathematicians. The minimal model program remains one of the most active areas
of research in algebraic geometry, with applications ranging from string theory
to number theory.

The K\"ahler--Einstein metric on Fano manifolds, now known to exist if and only
if the manifold is K-polystable, connects Mori's program to the theory of moduli
spaces. The moduli space of K-polystable Fano manifolds is a natural
compactification of the moduli space of smooth Fano manifolds, and its
construction is an active area of research.

The recent work of Birkar on the existence of minimal models for varieties of
general type uses the theory of asymptotic singularities and the LMMP (log
minimal model program). The key innovation is the use of the ``accidentally''
good behavior of log canonical singularities under the MMP, which ensures
termination of flips in many important cases.

The modern approach to the MMP uses the language of \textbf{singularities of pairs}.
A pair $(X, \Delta)$ consists of a normal variety $X$ and a $\mathbb{Q}$-divisor
$\Delta$ such that $K_X + \Delta$ is $\mathbb{Q}$-Cartier. The singularities of
the pair are classified by the discrepancies $a(E, X, \Delta)$, where $E$ is a
divisor over $X$. The key singularities are:
\begin{itemize}
    \item \textbf{Terminal}: $a(E, X, \Delta) > 0$ for all exceptional divisors $E$;
    \item \textbf{Canonical}: $a(E, X, \Delta) \geq 0$ for all exceptional divisors $E$;
    \item \textbf{Log terminal}: $a(E, X, \Delta) > -1$ for all exceptional divisors $E$;
    \item \textbf{Log canonical}: $a(E, X, \Delta) \geq -1$ for all exceptional divisors $E$.
\end{itemize}

The MMP is carried out in the category of pairs with log terminal singularities.
The key theorems---existence of flips, termination of flips, and existence of
minimal models---have all been proved for pairs with log terminal singularities
in all dimensions, thanks to the work of Birkar, Cascini, Hacon, and McKernan.

The theory of \textbf{asympotic singularities}, developed by Hacon and McKernan,
provides a powerful tool for proving the existence of flips in all dimensions.
The key idea is to study the asymptotic behavior of the linear series $|m(K_X +
\Delta)|$ for large $m$, and to use this to construct the flip as a limit of
a sequence of blow-ups.

The \textbf{ACC (ascending chain condition) conjecture} for log canonical thresholds
is a fundamental open problem in the MMP. The conjecture states that the set of
log canonical thresholds of pairs $(X, \Delta)$ is an ACC set, i.e., it contains
no infinite strictly increasing sequences. This conjecture, if true, would have
many important consequences for the structure of singularities.

The relationship between the MMP and the theory of \textbf{moduli spaces} has
become increasingly important. The moduli space of canonically polarized
varieties, constructed by Koll\'ar, Mori, and others, provides a geometric
framework for understanding the MMP. The moduli space of stable varieties
generalizes the moduli space of curves and provides a compactification of the
moduli space of smooth varieties of general type.

Mori's work continues to inspire new research in algebraic geometry. The
development of the MMP for pairs with log terminal singularities, the study of
rational curves on Fano varieties, and the connections to mirror symmetry and
string theory are all active areas of research that build on Mori's
foundational contributions.

The \textbf{effective MMP} has emerged as a major research direction. The goal
is to make the abstract existence proofs of the MMP into effective algorithms
that can be used to construct explicit minimal models. Recent work by
Birkar, Cascini, Hacon, and McKernan has provided effective bounds on the
discriminant of the minimal model program, but the problem of computing minimal
models explicitly remains challenging.

The study of \textbf{Fano varieties in positive characteristic} is another active
area of research. Mori's techniques, developed in characteristic zero, require
modification in positive characteristic. The existence of flips and minimal
models in positive characteristic is still largely open, although significant
progress has been made by Birkar, Waldron, and others.

The connection between the MMP and \textbf{non-Archimedean geometry} has led to
new insights. The theory of Berkovich spaces and the study of degenerations of
varieties have provided a geometric framework for understanding the MMP. The
work of Boucksom, Hsiao, Guedj, and others has shown that the MMP can be
formulated in terms of the existence of canonical metrics on non-Archimedean
line bundles.

\begin{center}\rule{0.5\linewidth}{0.5pt}\end{center}

\section*{Summary}
\addcontentsline{toc}{section}{Summary}

Mori proved the cone theorem and developed the minimal model program for
three-folds, showing that algebraic varieties can be classified up to birational
equivalence by constructing canonical representatives. His bend-and-break
technique and the theory of extremal rays provide the foundation for modern
birational geometry, connecting to K\"ahler geometry, Fano manifolds, and the
classification of rational curves. The classification of Fano three-folds
(105 deformation families) and the proof of flip existence in dimension 3 are
among the greatest achievements in algebraic geometry. The minimal model program,
now extended to all dimensions, remains one of the most active areas of
research in mathematics.
